\section{The exact numerical calculus retained by the construction}
\label{sec:logical-basis}

The physical optimization preserves a particular classical numerical
calculus, rather than replacing PA by an arbitrary small checker.
Its codes share the natural numbers with variable labels. Let
\[
 p(a,b)=\frac{(a+b)(a+b+1)}2+b,
 \qquad I(A,B)=4p(A,B),
\]
and define membership, universal quantification, and equality by
\[
 M(a,b)=4p(a,b)+1,\qquad
 U(a,B)=4p(B,a)+2,\qquad E(a,b)=4p(a,b)+3.
\]
In formulas below, an arrow denotes the constructor $I$. Code zero is
truth, and code one is the false code $M(0,0)$. Put
$N(B)=I(B,1)$, $B\wedge C=N(I(B,N(C)))$, and
$\exists x_a B=N(U(a,N(B)))$.
The quantifier's arguments are transposed in its numerical encoding.
The exception $I(0,0)=0$ is handled explicitly: zero is decoded as truth,
and the proof translations use the fixed equivalence between truth and
truth implying truth. No argument assumes that all constructor images
are syntactically disjoint at zero.

\subsection{Axiom heads and gate}

The six logical heads, for every $A,D,H\in\mathbb N$ and in the notation
used by the actual emitter, are
\begin{align*}
L_1&=(\exists x_D\,U(D,H))\to H,\\
L_2&=U(D,U(A,H))\to U(A,U(D,H)),\\
L_3&=M(H,D)\to(E(H,A)\to M(A,D)),\\
L_4&=E(H,D)\to(M(A,H)\to M(A,D)),\\
L_5&=\exists x_D E(H,D),\\
L_6&=U(D,A)\to\bigl(U(D,A\to H)\to U(D,H)\bigr).
\end{align*}
The implicational head ranges over $D,A,H,d\in\mathbb N$ with $d\le D$:
\[
 \rho(D,A,H,d)=((H\to d)\to D)\to((D\to H)\to(A\to H)),
 \qquad d\le D.
\]
The separate diagonal head, for every $B\in\mathbb N$, is
\[
                    \Gamma(B)=B\to U(B,B).
\]
For the guarded head, for every $a,b,c\in\mathbb N$, define
\begin{align*}
 V(a,b,c)&=M(a,b)\to U(c,M(a,b)),\\
 P(a,b,c)&=c\to M(c,b),\\
 S(a,b,c)&=U\bigl(b,(M(a,b)\to P(a,b,c))\wedge(M(c,b)\to c)\bigr),\\
 Q(a,b,c)&=U(c,S(a,b,c)\to N(P(a,b,c))),\\
 G(a,b,c)&=N(V(a,b,c)\to Q(a,b,c)),\\
 K(a,b,c)&=\begin{cases}G(a,b,c),&c\ne a\text{ and }c\ne b,\\0,&\text{otherwise.}\end{cases}
\end{align*}
The calculus has these axiom instances, modus ponens, and generalization
of theorems. The emitter additionally pushes a generalization of its old
cached theorem, as specified in the next section; that entry is not a new
axiom schema. The supplied development of the restricted implicational
basis recovers the required classical propositional interfaces. Its
proof-producing translations are part of the inherited calibration.

\subsection{Freshness and the separation of the diagonal head}

Every variable label occurring in a nonzero formula code $B$ is strictly
less than $B$. Induction on decoding proves this from the strict component
bounds of nonzero constructors; truth has no variables. Consequently
$x_B$ is absent from the formula with code $B$, and $\Gamma(B)$ is an
unconditional instance of vacuous generalization. It is safe without the
two comparisons guarding $K$.

The earlier bundled-diagonal gate was
\[
 F_\Gamma(a,b,c)=N\bigl(\Gamma(c)\to(V(a,b,c)\to Q(a,b,c))\bigr).
\]
These translations use the recovered classical calculus, including its
derived ex-falso principle. To see the latter explicitly, put
$E_B=\exists x_0 U(0,B)=N(U(0,N(U(0,B))))$. The recovered propositional
K law gives $1\to E_B$, while $L_1$ at $D=0,H=B$ gives $E_B\to B$;
implication composition therefore gives $1\to B$. With this principle,
the classical templates derive both $\Gamma(c)$ and $G(a,b,c)$
from the bundled formula and conversely. Every desired diagonal instance can be
obtained from a guarded bundled instance with $a=b=c+1$. Thus the old
bundled and current separate calculi have primitive-recursive proof
translations preserving conclusions. At command level the separation
adds a saved stack entry, which is why the correct ten-entry selectors
must be proved instead of copied from the old nine-entry interface.

The guards on $K$ cannot be removed merely because a separation subformula
looks easier when labels coincide. The same gate asserts
$V(a,b,c)$. For $c=a$, its antecedent can hold for a finite-set code while
universal quantification over that set variable makes the consequent false;
the case $c=b$ has the analogous problem in the element position. The
present construction preserves both guards.

\subsection{Logical bridge used in the theorem}

The inherited logical calibration\cite{pa195,pa194,pa195c} identifies refutations of this precise calculus
with PA refutations through syntactic transformations in both directions.
On the PA-to-calculus side, diagonal vacuity recovers the equality and
membership interfaces needed by the inherited PA-refutation compiler.
On the reverse side, membership is interpreted as binary membership of a
finite-set code, equality as arithmetic equality, and quantification as
first-order arithmetic quantification. Finite comprehension is supplied by
PA induction, with a uniform primitive-recursive procedure producing the
proof instance for each input formula.

For definiteness the binary-membership formula may be taken as
\[
 \operatorname{Bit}(s,z)\equiv
 \exists u,v\,[s=(2u+1)2^z+v\ \wedge\ v<2^z].
\]
PA proves $\operatorname{Bit}(s,z)\to z<s$, hence the negation of the
translation of $M(0,0)$. It also proves, for every arithmetic formula
$\theta(z,\vec u)$,
\[
 \forall n\,\exists s<2^n\,\forall z<n\,
 (\operatorname{Bit}(s,z)\leftrightarrow\theta(z,\vec u)).
\]
The proof extends the previous code by either zero or $2^n$ according to
$\theta(n,\vec u)$. Instantiating this induction-proof pattern is a
syntactic operation on formula codes, not a uniform truth predicate.
When the two pointed-separation labels coincide, the inner binder changes
both occurrences; the reference proof treats that case separately.

The next section gives the complete changed-controller proof compiler and
states how it composes with these logical translations. The detailed
inherited diagonal and separation arguments are retained in
\file{reference/logical_calibration/} and \file{references/}; their exact
input identities are recorded. The new physical proof does not derive
the logical calibration merely from observed nonhalting.
